\chapter{牛顿类迭代与实现细节}

\section{残差与雅可比构造}
对每个观测点，先调用泰勒积分器得到预测值，再计算残差。雅可比可由两种方式获得：
\begin{enumerate}
  \item 解析/符号导数方式：精度高，但推导复杂；
  \item 差分近似方式：实现简单，但需控制扰动尺度。
\end{enumerate}
综合可维护性与稳定性，本文推荐“核心参数解析导数 + 其余参数差分修正”的混合策略。

\section{阻尼高斯牛顿策略}
为避免远离最优点时出现发散，采用阻尼更新：
\[
\left(\mathbf{J}^\mathsf{T}\mathbf{J}+\lambda \mathbf{I}\right)\Delta\boldsymbol{\theta}
=-\mathbf{J}^\mathsf{T}\mathbf{r}.
\]
当目标函数下降不明显时增大 $\lambda$，下降明显时减小 $\lambda$。

\section{停止准则}
设置三类停止条件：
\begin{itemize}
  \item 参数步长：$\|\Delta\boldsymbol{\theta}\|_2 < \varepsilon_\theta$；
  \item 梯度范数：$\|\mathbf{g}\|_2 < \varepsilon_g$；
  \item 最大迭代次数：$k\ge k_{\max}$。
\end{itemize}

\section{复杂度分析}
设状态维度 $n$、参数维度 $p$、观测点数 $m$。每次迭代主要成本包括：
\begin{enumerate}
  \item 积分求解：$O(m\cdot n \cdot q)$，$q$ 为导数级数计算常数；
  \item 构造雅可比：$O(mnp)$；
  \item 线性方程求解：$O(p^3)$。
\end{enumerate}
当 $p$ 远小于 $mn$ 时，积分与雅可比评估是主导开销。
